% Part: incompleteness
% Chapter: theories-computability
% Section: extensions-of-q-not-decidable

\documentclass[../../../include/open-logic-section]{subfiles}

\begin{document}

\olfileid{inc}{tcp}{cqn}

\olsection{د~$\Th{Q}$ سازګارې غځونې ناپرېکړه‌وړې دي}

\begin{explain}
په ياد ولرئ چې يوه تيوري هله~\emph{سازګاره} ده چې د هېڅ
فارمول~$!A$ دپاره هم~$!A$ او هم~$\lnot !A$ دواړه ثابت نۀ کړي۔ ځکه
له تناقض څخه هر څه لازمېږي، يوه ناسازګاره تيوري ابتذالي وي: هره
جمله پکې ثابتېدونکې وي۔ څرګنده ده چې که يوه تيوري~$\omega$-سازګاره
وي، نو سازګاره هم ده۔ خو يوازې سازګاره اوسېدل کمزورے شرط دے (يعنې
داسې تيورۍ شته چې سازګارې دي خو~$\omega$-سازګارې نۀ دي)۔ په
\olref[oqn]{thm:oconsis-q} کښې فرضيه ساده سازګارۍ ته کمزورې کولے
شو، او په دې توګه يوه پياوړې قضيه ترلاسه کوو۔
\end{explain}

\begin{lem}
هېڅ «نړيواله محاسبه کېدونکې اړيکه» نشته۔ يعنې داسې دوه‌ځاييزه
محاسبه کېدونکې اړيکه~$R(x,y)$ نشته چې لاندې خاصيت ولري: هر کله
$S(y)$ يوه يوځاييزه محاسبه کېدونکې اړيکه وي، داسې~$k$ وي چې د هر
$y$ دپاره~$S(y)$ هله او يوازې هله رښتيا وي چې~$R(k,y)$ رښتيا وي۔
\end{lem}

\begin{proof}
فرض کړئ~$R(x,y)$ يوه نړيواله محاسبه کېدونکې اړيکه وي۔ پرېږدئ
$S(y)$، اړيکه~$\lnot R(y,y)$ وي۔ ځکه~$S(y)$ محاسبه کېدونکې ده، د
کوم~$k$ دپاره~$S(y)$ له~$R(k,y)$ سره هم‌ارزه ده۔ خو بيا~$S(k)$ هم
له~$R(k,k)$ او هم له~$\lnot R(k,k)$ سره هم‌ارزه ده، چې تناقض دے۔
\end{proof}


\begin{thm}
پرېږدئ~$\Th{T}$ هره هغه سازګاره تيوري وي چې~$\Th{Q}$ پکې شامله
وي۔ بيا~$\Th{T}$ د پرېکړې وړ نۀ ده۔
\end{thm}


\begin{proof}
فرض کړئ~$\Th{T}$ د~$\Th{Q}$ يوه سازګاره او د پرېکړې وړ غځونه وي۔
موږ به د~$\Th{T}$ په کارولو د يوې نړيوالې محاسبه کېدونکې اړيکې په
تعريفولو سره تناقض ترلاسه کړو۔

پرېږدئ~$R(x,y)$ هله او يوازې هله صدق وکړي چې
\begin{quote}
$x$ د~$!D(u)$ فارمول کوډوي، او~$\Th{T}$، $!D(\num y)$ ثابتوي۔
\end{quote}
ځکه چې فرض کوو~$\Th{T}$ د پرېکړې وړ ده، $R$ محاسبه کېدونکې ده۔
راځئ وښيو چې~$R$ نړيواله ده۔ که~$S(y)$ هره محاسبه کېدونکې اړيکه
وي، نو په~$\Th{Q}$ (او له دې کبله په~$\Th{T}$) کښې د کوم
فارمول~$!D_S(u)$ په وسيله تمثيلېږي۔ بيا د هر~$n$ دپاره لرو:
\begin{eqnarray*}
S(n) & \lif & T \vdash !D_S(\num n) \\
& \lif & R(\Gn{!D_S(u)}, n)
\end{eqnarray*}
او
\begin{eqnarray*}
\lnot S(n) & \lif & T \vdash \lnot !D_S(\num n) \\
& \lif & T \not\vdash !D_S(\num n) \quad \text{(ځکه~$\Th{T}$ سازګاره ده)} \\
& \lif & \lnot R(\Gn{!D_S(u)}, n).
\end{eqnarray*}
\textbf{د ژباړې د نسخې سمون (OLCMP-092):} سرچينې په دواړو
ماوراءژبنیو اړيکو کښې طبيعي عدد د څيز-ژبې د عددنښې په بڼه ليکلے و؛
دلته اړيکه پر طبيعي عدد تطبيق شوې، او عددنښه يوازې د تمثيلوونکي
فارمول دننه پاتې ده۔
يعنې د هر~$y$ دپاره~$S(y)$ هله او يوازې هله رښتيا وي چې
$R(\Gn{!D_S(u)}, y)$ رښتيا وي۔ نو~$R$ نړيواله ده، او هغه تناقض مو
ترلاسه کړ چې غوښت۔
\end{proof}

پرېږدئ «رښتينی حساب» تيوري~$\Setabs{!A}{ \Sat{\Nat}{!A}}$ وي؛
يعنې د حساب د ژبې د هغو جملو سټ چې په معياري تفسير کښې رښتيا دي۔

\begin{cor}
رښتينی حساب د پرېکړې وړ نۀ دے۔
\end{cor}

%In Section~\olref{undefinability:truth:section} we will state a stronger
%result, due to Tarski.

\end{document}
